
\chapter{2022年上海卷（春）}

\section{填空题}

\begin{problem}
已知 \(z=2+i\)，则 \(\bar z=\fillinblank{}\).

\end{problem}

\begin{answer}
\(2-i\)
\end{answer}

\begin{solution}
由共轭复数定义，\(\bar z=2-i\).
\end{solution}

\begin{problem}
已知 \(A=(-1,2)\)，\(B=(1,3)\)，则 \(A\cap B=\fillinblank{}\).

\end{problem}

\begin{answer}
\((1,2)\)
\end{answer}

\begin{solution}
由区间交集的定义可知，\(A\cap B=(1,2)\).
\end{solution}

\begin{problem}
不等式 \(\frac{x-1}{x}<0\) 的解集为\fillinblank{}.

\end{problem}

\begin{answer}
\(\{x\mid 0<x<1\}\)
\end{answer}

\begin{solution}
分式不等式为负，需分子分母异号，解得 \(0<x<1\).
\end{solution}

\begin{problem}
已知 \(\tan\alpha=3\)，则 \(\tan\!\left(\alpha+\frac\pi4\right)=\fillinblank{}\).

\end{problem}

\begin{answer}
\(-2\)
\end{answer}

\begin{solution}
由和角公式，\(\tan\!\left(\alpha+\frac\pi4\right)=\frac{\tan\alpha+1}{1-\tan\alpha}=\frac{3+1}{1-3}=-2\).
\end{solution}

\begin{problem}
已知方程组 \(\begin{cases}x+my=2,\\mx+16y=8,\end{cases}\) 有无穷解，则 \(m\) 的值为\fillinblank{}.

\end{problem}

\begin{answer}
\(4\)
\end{answer}

\begin{solution}
方程组有无穷解，说明两方程对应直线重合，故系数成比例. 由 \(\frac{1}{m}=\frac{m}{16}=\frac{2}{8}=\frac14\) 可得 \(m=4\).
\end{solution}

\begin{problem}
已知函数 \(f(x)=x^3\) 的反函数为 \(y=f^{-1}(x)\)，则 \(f^{-1}(27)=\fillinblank{}\).

\end{problem}

\begin{answer}
\(3\)
\end{answer}

\begin{solution}
由 \(y=x^3\) 得 \(x=\sqrt[3]{y}\)，故 \(f^{-1}(x)=\sqrt[3]{x}\)，从而 \(f^{-1}(27)=3\).
\end{solution}

\begin{problem}
在 \(\left(x^3+\frac1x\right)^{12}\) 的展开式中，含 \(\frac1{x^4}\) 项的系数为\fillinblank{}.

\end{problem}

\begin{answer}
\(66\)
\end{answer}

\begin{solution}
通项为 \(T_{r+1}=C_{12}^r(x^3)^{12-r}\left(\frac1x\right)^r=C_{12}^r x^{36-4r}\). 令 \(36-4r=-4\)，得 \(r=10\). 故所求系数为 \(C_{12}^{10}=66\).
\end{solution}

\begin{problem}
在 \(\triangle ABC\) 中，\(\angle A=\frac\pi3\)，\(AB=2\)，\(AC=3\)，则 \(\triangle ABC\) 的外接圆半径为\fillinblank{}.

\end{problem}

\begin{answer}
\(\frac{\sqrt{21}}{3}\)
\end{answer}

\begin{solution}
由余弦定理，\(BC^2=AB^2+AC^2-2\cdot AB\cdot AC\cos A=4+9-12\cdot\frac12=7\)，
故 \(BC=\sqrt7\). 再由正弦定理得 \(2R=\frac{BC}{\sin A}=\frac{\sqrt7}{\sqrt3/2}\)，故 \(R=\frac{\sqrt{21}}{3}\).
\end{solution}

\begin{problem}
已知有 \(1,2,3,4\) 四个数字组成无重复数字，则比 \(2134\) 大的四位数的个数为\fillinblank{}.

\end{problem}

\begin{answer}
\(17\)
\end{answer}

\begin{solution}
按千位分类，千位为 \(3,4\) 时，共 \(2A_3^3=12\) 个；千位为 \(2\) 且百位为 \(3\) 或 \(4\) 时，共 \(2A_2^2=4\) 个；千位为 \(2\) 且百位为 \(1\) 时，只有 \(2143\) 比 \(2134\) 大，贡献 \(1\) 个，合计 \(17\) 个.
\end{solution}

\begin{problem}
在 \(\triangle ABC\) 中，\(\angle C=\frac\pi2\)，\(AC=BC=2\)，\(M\) 为 \(AC\) 的中点，\(P\) 在线段 \(AB\) 上，则 \(\overrightarrow{MP}\cdot\overrightarrow{CP}\) 的最小值为\fillinblank{}.

\end{problem}

\begin{answer}
\(\frac78\)
\end{answer}

\begin{solution}
建立平面直角坐标系，取 \(C(0,0)\)，\(A(2,0)\)，\(B(0,2)\)，则 \(M(1,0)\). 设 \(P(x,2-x)\)（\(0\le x\le2\)），则 \(\overrightarrow{MP}=(x-1,2-x)\)，\(\overrightarrow{CP}=(x,2-x)\). 于是 \(\overrightarrow{MP}\cdot\overrightarrow{CP}=x(x-1)+(2-x)^2=2x^2-5x+4\). 在 \(x\in[0,2]\) 上其最小值为 \(\frac78\).
\end{solution}

\begin{problem}
已知双曲线 \(\frac{x^2}{a^2}-y^2=1\,(a>0)\)，双曲线上右支上有任意两点 \(P_1(x_1,y_1),P_2(x_2,y_2)\)，满足 \(x_1x_2-y_1y_2>0\) 恒成立，则 \(a\) 的取值范围是\fillinblank{}.

\end{problem}

\begin{answer}
\(a\ge1\)
\end{answer}

\begin{solution}
设右支上点 \(P(x,y)\) 与渐近线 \(y=\pm \frac{x}{a}\) 的夹角分别记为 \(\theta\). 条件 \(x_1x_2-y_1y_2>0\) 等价于两向量夹角小于 \(\frac\pi2\). 因而渐近线张开的角不超过 \(\frac\pi2\)，即 \(2\arctan\frac1a\le\frac\pi2\). 故 \(\frac1a\le1\)，即 \(a\ge1\).
\end{solution}

\begin{problem}
已知 \(f(x)\) 为奇函数，当 \(x\in(0,1]\) 时，\(f(x)=\ln x\)，且 \(f(x)\) 关于直线 \(x=1\) 对称，设 \(f(x)=x+1\) 的正数解依次为 \(x_1,x_2,x_3,\ldots,x_n,\ldots\)，则 \(\lim\limits_{n\to\infty}(x_{n+1}-x_n)=\fillinblank{}\).

\end{problem}

\begin{answer}
\(2\)
\end{answer}

\begin{solution}
由奇函数与关于直线 \(x=1\) 对称可推出 \(f(x)\) 以 \(4\) 为周期. 作图可知 \(f(x)=x+1\) 的正数解在相邻周期内间隔恒为 \(2\)，故极限为 \(2\).
\end{solution}
\section{单选题}

\begin{problem}
下列函数定义域为 \(\R\) 的是（\quad）

\choices
    {\(y=x^{-2}\)}
    {\(y=x^{-1}\)}
    {\(y=x^{\frac13}\)}
    {\(y=x^{\frac12}\)}

\end{problem}

\begin{answer}
C
\end{answer}

\begin{solution}
\(x^{1/3}\) 在全体实数上有定义，其余选项都有限制. 故选 C.
\end{solution}

\begin{problem}
已知 \(a>b>c>d\)，下列选项中正确的是（\quad）

\choices
    {\(a+d>b+c\)}
    {\(a+c>b+d\)}
    {\(ad>bc\)}
    {\(ac>bd\)}

\end{problem}

\begin{answer}
B
\end{answer}

\begin{solution}
由 \(a> b\) 与 \(c>d\) 相加，得 \(a+c>b+d\). 故选 B.
\end{solution}

\begin{problem}
如图，上海海关大楼的上面可以看作一个正四棱柱，四个侧面有四个时钟，则相邻两个时钟的时针从 \(0\) 时转到 \(12\) 时（含 \(0\) 时不含 \(12\) 时）的过程中，能够相互垂直（\quad）次.

\choices
    {\(0\)}
    {\(2\)}
    {\(4\)}
    {\(12\)}

\bitmapfigure[max width=0.22\linewidth,max totalheight=4.8cm]{img/2022/shanghai_spring/q15_fig1.png}

\end{problem}

\begin{answer}
B
\end{answer}

\begin{solution}
相邻两个钟面上的时针只有在 \(3\) 时和 \(9\) 时的位置才互相垂直，其余时刻不满足，故共 \(2\) 次.
\end{solution}

\begin{problem}
已知等比数列 \(\{a_n\}\) 的前 \(n\) 项和为 \(S_n\)，前 \(n\) 项积为 \(T_n\)，则下列选项判断正确的是（\quad）

\choices
    {若 \(S_{2022}>S_{2021}\)，则数列 \(\{a_n\}\) 是递增数列}
    {若 \(T_{2022}>T_{2021}\)，则数列 \(\{a_n\}\) 是递增数列}
    {若数列 \(\{S_n\}\) 是递增数列，则 \(a_{2022}\ge a_{2021}\)}
    {若数列 \(\{T_n\}\) 是递增数列，则 \(a_{2022}\ge a_{2021}\)}

\end{problem}

\begin{answer}
D
\end{answer}

\begin{solution}
前 \(n\) 项积递增意味着每一步乘上的公比满足 \(q\ge1\)，从而 \(a_{2022}\ge a_{2021}\). 其余三项均可构造反例. 故选 D.
\end{solution}
\section{解答题}

\begin{problem}
如图，在圆柱 \(OO_1\) 中，底面半径为 \(1\)，\(AA_1\) 为圆柱 \(OO_1\) 的母线.

\begin{enumerate}
\item 若 \(AA_1=4\)，\(M\) 为 \(AA_1\) 的中点，求直线 \(MO_1\) 与底面的夹角大小；
\item 若圆柱的轴截面为正方形，求该圆柱的侧面积和体积.
\end{enumerate}

\bitmapfigure[max width=0.33\linewidth,max totalheight=4.8cm]{img/2022/shanghai_spring/q17_fig1.png}

\end{problem}

\begin{solution}
（1）直线 \(MO_1\) 与底面的夹角等于 \(\angle MO_1A_1\). 在直角三角形 \(MO_1A_1\) 中，\(MA_1=\frac12AA_1=2\)，\(O_1A_1=1\)，故 \(\tan\angle MO_1A_1=\frac{MA_1}{O_1A_1}=2\). 所以夹角为 \(\arctan 2\).

（2）轴截面为正方形，故 \(AA_1=2\). 于是侧面积为 \(2\pi r h=2\pi\cdot1\cdot2=4\pi\)，体积为 \(\pi r^2h=\pi\cdot1^2\cdot2=2\pi\).
\end{solution}

\begin{problem}
已知数列 \(\{a_n\}\)，\(a_2=1\)，\(\{a_n\}\) 的前 \(n\) 项和为 \(S_n\).

\begin{enumerate}
\item 若 \(\{a_n\}\) 为等比数列，\(S_2=3\)，求 \(\lim\limits_{n\to\infty}S_n\)；
\item 若 \(\{a_n\}\) 为等差数列，公差为 \(d\)，对任意 \(n\in\N^*\)，均满足 \(S_{2n}\ge n\)，求 \(d\) 的取值范围.
\end{enumerate}

\end{problem}

\begin{solution}
（1）由 \(S_2=a_1+a_2=3\)，且 \(a_2=1\)，得 \(a_1=2\). 等比数列公比 \(q=\frac12\). 故 \(S_n=\frac{2(1-(\frac12)^n)}{1-\frac12}=4\Bigl(1-\bigl(\tfrac12\bigr)^n\Bigr)\)，因此 \(\lim\limits_{n\to\infty}S_n=4\).

（2）由 \(a_2=1\) 得 \(a_1=1-d\). 又 \(S_{2n}=\frac{2n(a_1+a_{2n})}{2}=n(a_2+a_{2n-1})\)，故条件等价于 \(a_2+a_{2n-1}\ge1\)，即 \(1+\bigl(1+(2n-3)d\bigr)\ge1\). 分 \(n=1\) 与 \(n\ge2\) 讨论，可得 \(d\in[0,1]\).
\end{solution}

\begin{problem}
如图，矩形 \(ABCD\) 区域内，\(D\) 处有一棵古树，为保护古树，以 \(D\) 为圆心，\(DA\) 为半径划定圆 \(D\) 作为保护区域，已知 \(AB=30\)m，\(AD=15\)m，点 \(E\) 为 \(AB\) 上的动点，点 \(F\) 为 \(CD\) 上的动点，满足 \(EF\) 与圆 \(D\) 相切.

\begin{enumerate}
\item 若 \(\angle ADE=20^\circ\)，求 \(EF\) 的长；
\item 当点 \(E\) 在 \(AB\) 的什么位置时，梯形 \(FEBC\) 的面积有最大值，最大面积为多少？
\end{enumerate}

\end{problem}

\begin{solution}
（1）设 \(EF\) 与圆 \(D\) 相切于 \(H\)，连接 \(DH\). 因为 \(DH\perp EF\)，且 \(DH=AD=15\). 在直角三角形 \(HED\) 与 \(FHD\) 中，\(EH=15\tan20^\circ\)，\(HF=15\tan50^\circ\). 所以 \(EF=EH+HF=15(\tan20^\circ+\tan50^\circ)\approx23.3\text{ m}\).

（2）设 \(\angle ADE=\theta\). 则 \(AE=15\tan\theta\)，\(FH=15\tan(90^\circ-2\theta)\). 梯形 \(AEFD\) 的面积为 \(S=\frac{225}{4}\left(3\tan\theta+\frac1{\tan\theta}\right)\ge\frac{225\sqrt3}{2}\)，当 \(\tan\theta=\frac{\sqrt3}{3}\) 时取等号，此时 \(AE=15\tan\theta=5\sqrt3\approx8.7\text{ m}\). 故当 \(AE=8.7\) m 时，梯形 \(FEBC\) 的面积最大，最大值为 \(15\times30-\frac{225\sqrt3}{2}\approx255.14\text{ m}^2\).
\end{solution}

\begin{problem}
在椭圆 \(\Gamma:\frac{x^2}{a^2}+y^2=1\) 中，直线 \(l:x=a\) 上有两点 \(C,D\)（\(C\) 点在第一象限），左顶点为 \(A\)，下顶点为 \(B\)，右焦点为 \(F\).

\begin{enumerate}
\item 若 \(\angle AFB=\frac\pi6\)，求椭圆 \(\Gamma\) 的标准方程；
\item 若点 \(C\) 的纵坐标为 \(2\)，点 \(D\) 的纵坐标为 \(1\)，则 \(BC\) 与 \(AD\) 的交点是否在椭圆上？请说明理由；
\item 已知直线 \(BC\) 与椭圆 \(\Gamma\) 相交于点 \(P\)，直线 \(AD\) 与椭圆 \(\Gamma\) 相交于点 \(Q\)，若 \(P\) 与 \(Q\) 关于原点对称，求 \(\lvert CD\rvert\) 的最小值.
\end{enumerate}

\end{problem}

\begin{solution}
（1）设椭圆焦距为 \(c\)，则 \(A(-a,0)\)，\(B(0,-1)\)，\(F(c,0)\). 由 \(\tan\angle AFB=\frac{1}{c}=\tan\frac\pi6=\frac{\sqrt3}{3}\)，得 \(c=\sqrt3\). 又 \(a^2=c^2+1=4\)，故椭圆方程为 \(\frac{x^2}{4}+y^2=1\).

（2）由 \(B(0,-1),C(a,2)\) 得 \(BC:\ y=\frac3a x+1\)；由 \(A(-a,0),D(a,1)\) 得 \(AD:\ y=\frac1{2a}(x+a)\). 联立得交点为 \(\left(\frac{3a}{5},\frac45\right)\). 代入椭圆方程得 \(\left(\frac{3a}{5}\right)^2\!\!\bigg/\!a^2+\left(\frac45\right)^2=\frac{9}{25}+\frac{16}{25}=1\)，故交点在椭圆上.

（3）设 \(P(a\cos\theta,\sin\theta)\)，\(Q(-a\cos\theta,-\sin\theta)\). 由直线方程可表示 \(C,D\) 的坐标，并将 \(\lvert CD\rvert\) 化为关于 \(t=\tan\frac\theta2\) 的式子，可得 \(\lvert CD\rvert=2\left(\frac1{1-t}+\frac1t-1\right)\ge6\). 故最小值为 \(6\).
\end{solution}

\begin{problem}
已知 \(x_1,x_2\in\R\)，若对任意的 \(x_2-x_1\in S\)，\(f(x_2)-f(x_1)\in S\)，则称 \(f(x)\) 是在 \(S\) 关联的.

\begin{enumerate}
\item 判断和证明 \(f(x)=2x-1\) 是否在 \([0,+\infty)\) 关联？是否在 \([0,1]\) 关联？
\item 若 \(f(x)\) 是在 \(\{3\}\) 关联的，且 \(f(x)\) 在 \(x\in[0,3)\) 时 \(f(x)=x^2-2x\)，求解不等式 \(2\le f(x)\le3\)；
\item 证明：\(f(x)\) 是 \(\{1\}\) 关联的，且是在 \([0,+\infty)\) 关联的，当且仅当“\(f(x)\) 在 \([1,2]\) 是关联的”.
\end{enumerate}

\end{problem}

\begin{solution}
（1）\(f(x)=2x-1\) 在 \([0,+\infty)\) 关联，但在 \([0,1]\) 不关联. 若 \(x_2-x_1\in[0,+\infty)\)，则 \(f(x_2)-f(x_1)=2(x_2-x_1)\in[0,+\infty)\)，故在 \([0,+\infty)\) 关联. 取 \(x_1=1,x_2=0\)，则 \(x_2-x_1=-1\notin[0,1]\) 不在定义范围内；再取 \(x_1=1,x_2=2\)，则差为 \(1\in[0,1]\)，而 \(f(2)-f(1)=2\notin[0,1]\)，故在 \([0,1]\) 不关联.

（2）由“在 \(\{3\}\) 关联”可知 \(f(x+3)-f(x)=3\). 先由 \(x\in[0,3)\) 时 \(f(x)=x^2-2x\) 求得值域，再用平移关系得到 \(x\in[3,6)\) 的值域. 逐段解不等式可得 \(2\le f(x)\le3\) 的解集为 \([\sqrt3+1,5]\).

（3）若同时满足“在 \(\{1\}\) 关联”且“在 \([0,+\infty)\) 关联”，则对任意 \(1\le x_2-x_1\le2\)，有 \(1\le f(x_2)-f(x_1)\le2\)，故在 \([1,2]\) 关联. 反过来，若在 \([1,2]\) 关联，则可推出 \(f(x+1)-f(x)=1\)，再由任意正差分解到整数部分与 \([1,2]\) 部分，可知在 \(\{1\}\) 关联且在 \([0,+\infty)\) 关联. 综上，结论成立.
\end{solution}
